%! TeX program = lualatex
\documentclass[10pt,a4paper,sans]{article} % Font sizes: 10, 11, or 12; paper sizes: a4paper, letterpaper, a5paper, legalpaper, executivepaper or landscape; font families: sans or roman
\usepackage[hidelinks]{hyperref}

\usepackage[export]{adjustbox}

\usepackage{float}
\usepackage{array}
\usepackage{multicol}
\usepackage{multirow}
\usepackage[most]{tcolorbox}

\usepackage[scale=1]{geometry} % Reduce document margins
%\setlength{\hintscolumnwidth}{3.5cm} 
%\newcommand{\trig}{$\blacktriangleright$\hspace{1mm}}
\usepackage{color}
\usepackage{xcolor}
\usepackage{amssymb}
\usepackage{lipsum} % Used for inserting dummy 'Lorem ipsum' text into the template
\usepackage{fontawesome5} %funny symbols
\usepackage{fontspec} %fonts: require LuaTex or XeTex
\setmainfont{Latin Modern Roman}
\setsansfont{Latin Modern Sans}
\renewcommand{\familydefault}{\sfdefault} % Set sans as the default
\usepackage[fontsize=10.85pt]{fontsize} %fuck yeah I can use floats as fontsize
\usepackage{setspace}
\setstretch{1.0} % Force tight single spacing
\setlength{\parskip}{0pt} % Remove space between paragraphs
\setlength{\parsep}{0pt}
\setlength{\parindent}{0pt}
\renewcommand{\baselinestretch}{1.0} % moderncv uses a tight stretch
\setlength{\lineskip}{0pt}
\setlength{\normallineskip}{0pt}

\usepackage{enumitem}
\setlist[itemize]{nosep, topsep=0pt, partopsep=0pt, leftmargin=0.8em}

\setlength{\tabcolsep}{0pt} % Remove default padding between columns
\renewcommand{\arraystretch}{1.0} % Ensure rows aren't stretched

\usepackage{afterpage} %page color i guess
\usepackage{tikz} %its tikz, dummy
\usepackage{todonotes}
\usepackage{framed} %frames
\usepackage{svg}
\usepackage{circledsteps}

%colors

\definecolor{green1}{rgb}{0.35,0.70,0.30}% green
\definecolor{darkgrey}{rgb}{0.529,0.52,0.588}% dark grey
\definecolor{bluegrey}{rgb}{0.529,0.52,0.588}% blue gray
\definecolor{bluelightgrey}{rgb}{0.935,0.925,1}% blue gray

%\definecolor{darkdarkgrey}{rgb}{0.184,0.184,0.227}% dark dark grey

\definecolor{darkdarkgrey}{rgb}{0.22,0.22,0.23}% dark dark grey
\definecolor{draculared}{rgb}{1, 0.33, 0.33}% dracula red



\definecolor{color0}{rgb}{0,0,0}
\definecolor{color1}{rgb}{0.529,0.52,0.588}% blue gray
\definecolor{color2}{rgb}{0.45,0.45,0.45}% dark grey

\definecolor{dockercolor}{HTML}{384d54}
\definecolor{cppcolor}{HTML}{f34b7d}
\definecolor{reponamecolor}{HTML}{4493f8}

%override commands

\newlength{\hintscolumnwidth}
\setlength{\hintscolumnwidth}{1cm}

\renewcommand{\section}[1]{%
  \par\vspace{2.5ex}%
  \phantomsection{}%
  \begin{tabular}{@{}p{\hintscolumnwidth}@{\hspace{0.2cm}}p{0.8\textwidth}@{}}
    % This is the "bar" feature:
    \raggedleft\raisebox{0.2ex}{\color{bluegrey}\rule{\hintscolumnwidth}{0.95ex}} & 
    \Large\bfseries\upshape\textcolor{bluegrey}{#1}
  \end{tabular}
  \par\vspace{1ex}%
}

%New commands
\newcommand{\highlight}[1]{\textbf{\textcolor{draculared}{#1}}}
\newcommand{\CPP}{C\nolinebreak\hspace{-.05em}\raisebox{.4ex}{\tiny \textbf{+}}\nolinebreak\hspace{-.10em}\raisebox{.4ex}{\tiny \textbf{+}}\hspace{1mm}}

\newcommand{\Cs}{C\nolinebreak\raisebox{.4ex}{\scriptsize \textbf{\#}}}
%font stuff
%set font for my name
\setmonofont[Path = fonts/, Scale=1]{FredokaOne-Regular}
%defaut font
\defaultfontfeatures{LetterSpace=100.0}

\newcommand{\trig}{\faAngleRight\hspace{1mm}}
\newcommand{\skills}{\faAlignLeft\hspace{1mm}}
\newcommand{\academicitem}[4]{
  \begin{center}
    \begin{tabular}{p{0.14\textwidth} p{0.01\textwidth} !{\vrule width 2pt} p{0.015\textwidth} p{0.8\textwidth} } 
        \centering #1& & & \textbf{#2} \small\textcolor{darkgrey}{\textit{(#3)}} \newline\normalsize \trig \small #4 \normalsize\\ 
    \end{tabular}
  \end{center}
}

%date, title, address, subtitle
\newcommand{\academicitemafter}[4]{
    \begin{center}
        \begin{tabular}{p{0.14\textwidth} p{0.01\textwidth} !{\vrule width 2pt} p{0.015\textwidth} p{0.8\textwidth} } 
            \centering \small\textcolor{darkgrey}{\textit{(#3)}}\newline \normalsize#1& & & \textbf{#2}  \newline\normalsize \trig \small #4 \normalsize\\ 
        \end{tabular}
    \end{center}
}
% #1: id, #2: institution, #3: logo, #4: location, #5: start, #6: end, #7: desc, #8: skills
\newcommand{\academicitemtimeline}[8]{
    \noindent
    \begin{minipage}[t]{0.15\textwidth}
        \vspace{0pt}
        \centering
        \adjincludegraphics[width=0.4\linewidth, keepaspectratio, valign=t]{#3}
        \par\vspace{2mm}
        \footnotesize #5\par
        \footnotesize #6
        \vspace{2mm}
    \end{minipage}%
    \begin{minipage}[t]{0.82\textwidth}
        \vspace{0pt}
        \normalsize \textbf{#2} \small\textcolor{darkgrey}{\textit{(#4)}} \par
        \vspace{1mm}
        \footnotesize \trig #7
    \end{minipage}


    \vspace{1mm}
    \hspace{3mm}
    \begin{minipage}[t]{0.90\textwidth}
      \footnotesize \skills \textbf{Skills:} #8
    \end{minipage}

    \par\vspace{2mm} 
}


\newcommand{\jobitem}[5]{
    \begin{center}
        \begin{tabular}{p{0.14\textwidth} p{0.008\textwidth} !{\vrule width 2pt} p{0.015\textwidth} p{0.8\textwidth} } 
            \centering #1& & & \textbf{#2} \small\textcolor{darkgrey}{\textit{(#3)}} \newline  #4
            \vspace{1mm}\newline \trig #5 \normalsize\\
        \end{tabular}
    \end{center}
}

% #1: id, #2: company, #3: logo, #4: location, #5: role, #6: start, #7: end, #8: desc, #9: skills
\newcommand{\jobitemtimeline}[9]{
    \noindent
    \begin{minipage}[t]{0.02\textwidth}
        \vspace{0pt} 
        \centering
        \begin{tikzpicture}[remember picture, baseline=(#1.north)]
            \node[circle, fill=bluegrey, inner sep=1.2mm] (#1) {};
            \draw[line width=1mm , bluegrey] (#1.center) -- (0,-1);
        \end{tikzpicture}
    \end{minipage}%
    \begin{minipage}[t]{0.15\textwidth}
        \vspace{0pt}
        \centering
        \adjincludegraphics[width=0.45\linewidth, keepaspectratio, valign=t]{#3}
        \par\vspace{2mm}
        \footnotesize #6\par
        \footnotesize #7
    \end{minipage}%
    \begin{minipage}[t]{0.8\textwidth}
        \vspace{0pt}
        \normalsize \textbf{#5} \par
        \footnotesize #2 \textcolor{darkgrey}{\textit{(#4)}} \par
        \vspace{1mm}
        \footnotesize \trig #8
    \end{minipage}

    \par\vspace{2mm}
    \noindent
    \begin{minipage}[t]{0.9\textwidth}
        \footnotesize \skills \textbf{Skills:} #9
    \end{minipage}
    
    \par\vspace{3mm} % Fixed gap between job blocks
}

\newcommand{\connectjobs}[2]{
    \begin{tikzpicture}[remember picture, overlay]
        \draw[line width=1mm , bluegrey] (#1.center) -- (#2.center);
    \end{tikzpicture}
}

\newcommand{\infoitem}[2]{
    \begin{center}
        \begin{tabular}{ p{0.15\textwidth} p{0.8\textwidth} } 
            \centering#1&\small#2
        \end{tabular}
    \end{center}
    \vspace{-8.5mm}
    \normalsize
}


% All information in this block is optional, comment out any lines you don't need
\title{Curriculum Vitae}

\setlength{\hintscolumnwidth}{1cm} 

%multi-language settings
%
\newif\ifen
\newif\iffr
\newif\ifpt

\newcommand{\en}[2][]{\ifen#2\else#1\fi}
\newcommand{\fr}[2][]{\iffr#2\else#1\fi}
\newcommand{\pt}[2][]{\ifpt#2\else#1\fi}

%
\newcommand{\brazil}{\en{Brazil}\fr{Brésil}\pt{Brasil}}
\newcommand{\france}{\pt[France]{França}}

\newcommand{\january}{Jan}
\newcommand{\february}{\en{Feb}\fr{Fév}\pt{Fev}}
\newcommand{\mars}{Mar}
\newcommand{\april}{\en{Apr}\fr{Avr}\pt{Abr}}
\newcommand{\may}{\en[Mai]{May}}
\newcommand{\june}{\fr[Jun]{Juin}}
\newcommand{\july}{\fr[Jul]{Juil}}
\newcommand{\august}{\en{Aug}\fr{Août}\pt{Ago}}
\newcommand{\september}{\pt[Sep]{Set}}
\newcommand{\october}{\pt[Oct]{Out}}
\newcommand{\november}{Nov}
\newcommand{\december}{\en{Dec}\fr{Déc}\pt{Dez}}
\newcommand{\currently}{\en{Currently}\fr{Actuellement}\pt{Atualmente}}
%%%%
\entrue
% \frtrue
%\pttrue

\begin{document}

\pagecolor{bluegrey}

\tcbset{grow to right by=0.1cm,colframe=darkdarkgrey,colback=darkdarkgrey,arc=0mm}
\begin{tcolorbox}
    \begin{minipage}[t]{0.32\textwidth}
        
        \vspace*{-5mm}
        \begin{tcolorbox}[grow to left by=0.6cm,colback=darkdarkgrey,colframe=darkdarkgrey,arc=0mm]
        \color{white}
        
            \begin{tcolorbox}[colback=bluegrey,enlarge top by=-5mm,enlarge left by=-10mm,arc=4mm,sharp corners=north]
            \vspace{2mm}
            
            \end{tcolorbox}

            \vspace*{-12mm}
            
            \begin{center}
                \includegraphics[width=0.52\textwidth]{((parent_prefix))pictures/photo_cv_circ_20260225.png}
            \end{center}
            
            \vspace*{-4mm}
            \normalsize{\textcolor{white}{\texttt{ \textbf{Joe \en{Ferreira Scholtz }\fr{FERREIRA SCHOLTZ}\pt{Ferreira Scholtz}}\raisebox{0mm}{\footnotesize{ \en{(she/her)}\fr{(elle/d'elle)}\pt{ela/dela}}}}}}
         
            \vspace*{-2mm}
            
            \section{\textbf{\faAddressCard \hspace{1mm}\en{Contact}\fr{Contact}\pt{Contato}}}
            
            \vspace*{-4mm}

            \infoitem{\faEnvelope}{\footnotesize joefscholtz@gmail.com}
            \vspace{1mm}
            
            \infoitem{\faMobile*}{+55 51 99917-2042}
            \vspace{1mm}
            
            \infoitem{\faMapMarked*}{Canoas, \brazil}
            \vspace*{-7mm}
            
            \begin{center}
                \begin{tabular}{p{0.012\textwidth} p{0.478\textwidth}p{0.49\textwidth}} 
                     & \infoitem{\href{https://www.linkedin.com/in/joe-scholtz/}{\faLinkedin}}{\href{https://www.linkedin.com/in/joe-scholtz/}{Joe Scholtz}} & \infoitem{\href{https://github.com/joefscholtz/}{\faGithub}}{\href{https://github.com/joefscholtz}{joefscholtz}}
                \end{tabular}
            \end{center}

            \vspace*{-10mm}

            \setlength{\hintscolumnwidth}{1cm}
            \section{\textbf{\skills
            \en{Skills}\fr{Compétences}\pt{Competências}}}
            \setlength{\hintscolumnwidth}{0.5cm}
            \vspace*{-5mm}
            \section{\textbf{\faTerminal\hspace{1mm}
            \en[\textit{Software}]{Software}
            }}
            \vspace*{-3mm}

            \textcolor{bluegrey}{\hrulefill\newline}
            \hspace*{-1.5mm}\textcolor{bluegrey}{\textbf{\en{Relevant projects}\fr{Projets pertinents}\pt{Projetos relevantes}}}

            \vspace*{-2mm}
            

            \pgfkeys{/csteps/inner ysep=7pt}
            \pgfkeys{/csteps/inner xsep=7pt}
            \begin{tcolorbox}[left = 2mm,right = 2mm,colback=darkdarkgrey,colframe=bluegrey,coltext=white,arc=1mm]
            \small 
            \vspace*{-1mm}
                \begin{center}
                    \begin{tabular}{m{0.1\textwidth} m{0.35\textwidth} m{0.15\textwidth} m{0.41\textwidth} } 
                        \includesvg{((parent_prefix))pictures/git_repo.svg} & \textcolor{reponamecolor}{\textbf{\href{https://github.com/joefscholtz/forg_bot}{forg\_bot}}} &  \Circled{\footnotesize{Public}} & \text{ }
                    \end{tabular}
                \end{center}
                \vspace*{-2.8mm}
                \scriptsize{Navigation with Curiosity Rover inspired mobile robot in ROS 2 }\\
                \vspace*{-1.5mm}
                
                \footnotesize{\CircledTop[fill color=cppcolor,outer color=cppcolor]{} \CPP}
                \vspace*{-0.5mm}
            \end{tcolorbox}

            \vspace*{-3mm}

            \begin{tcolorbox}[left = 2mm,right = 2mm,colback=darkdarkgrey,colframe=bluegrey,coltext=white,arc=1mm]
            \small 
            \vspace*{-1mm}
                \begin{center}
                    \begin{tabular}{m{0.1\textwidth} m{0.40\textwidth} m{0.15\textwidth} m{0.36\textwidth} } 
                        \includesvg{((parent_prefix))pictures/git_repo.svg} & \textcolor{reponamecolor}{\textbf{\href{https://github.com/joefscholtz/ros2_docker}{ros2\_docker}}} &  \Circled{\footnotesize{Public}} & \text{ }
                    \end{tabular}
                \end{center}
                \vspace*{-2.8mm}
                \scriptsize{My workflow to run ROS 2 applications inside Docker, including base stage to deploy appli} \ldots\\
                \vspace*{-1.5mm}
                
                \footnotesize{\CircledTop[fill color=dockercolor,outer color=dockercolor]{} Dockerfile}
                \vspace*{-0.5mm}
            \end{tcolorbox}
            
            % \textcolor{bluegrey}{\hrulefill\newline}
            % \hspace*{-1.5mm}\textcolor{bluegrey}{\textbf{\en{OS}\fr{Systèmes d'exploitation}\pt{Sistemas operacionais}}}
            
            % \en{Windows and Linux.}
            % \fr{Windows et Linux.}
            % \pt{Windows e Linux.}
            \vspace*{-3mm}
            
            \textcolor{bluegrey}{\hrulefill\newline}
            \hspace*{-1.5mm}\textcolor{bluegrey}{\textbf{\en{Languages, tools and frameworks}\fr{Langages, outils et  \textit{frameworks}}\pt{Linguagens, ferramentas e \textit{frameworks}}}}
            %frameworks = environnements de développement
            \vspace*{1mm}

        
            \small 

            \begin{itemize}
                \item[\faAngleRight] C, \en{Embedded }C\fr{ Embarqué}\pt{ Embarcado}, \en{Modern }\CPP\fr{ moderne}\pt{ moderno}.
                \item[\faAngleRight] ROS \en{and}\fr{et}\pt{e} ROS 2:
                    \begin{itemize}
                        \item URDF, ros2\_control, SLAM, Nav2, Open-RMF.
                        % \item URDF, xacro, Rviz, SLAM, MoveIt, Nav2.
                    \end{itemize}
                \item[\faAngleRight] Python:
                    \begin{itemize}
                        \item NumPy, SciPy, Pandas, Matplotlib.
                    \end{itemize}
                % \item[\faAngleRight] IEC 61131-3
                %     \begin{itemize}
                %         \item  Instruction List, Ladder Diagram, Structured Text, Function Block Diagram, Sequential Flow Chart.
                %     \end{itemize}
                % \item[\faAngleRight] Web Development
                %     \begin{itemize}
                %         \item HTML, CSS, Javascript, PHP.
                %     \end{itemize}
                % \item[\faAngleRight] \Cs, MATLAB, Java, Git, SQL, \LaTeX.
                \item[\faAngleRight] MATLAB, Git, Docker, \LaTeX, bash, IEC 61131-3 (ST, IL/ASSEMBLY, LADDER, FBD, SFC).
                
                % \item[\faAngleRight] \en{Communication Protocols}\fr{Protocoles de communication}\pt{Protocolos de comunicação}:
                % \begin{itemize}
                %     \item LPWA, 6LOWPAN, LORA, MQTT, TCP, TLS/SSL, Wifi, Bluetooth/BLE, ZIGBEE, RS422, RS485, CAN.
                % \end{itemize}
            \end{itemize}
            
            \normalsize
            
            \vspace*{-1mm}
            
            \textcolor{bluegrey}{\hrulefill\newline}
            \hspace*{-1.5mm}\textcolor{bluegrey}{\textbf{\pt[Simulation]{Simulação}}}
            
            \small 
            
            \begin{itemize}
                \item[\faAngleRight] Gazebo, Isaac Sim, Simulink, Micro-Cap, PSIM, Proteus, LabView. 
            \end{itemize}
            
            \normalsize
            
            \vspace*{-1mm}
            
            \textcolor{bluegrey}{\hrulefill\newline}
            \hspace*{-1.5mm}\textcolor{bluegrey}{\textbf{\en{CAD and 3D modelling}\fr{CAO et modélisation 3D}\pt{\textit{CAD} e modelagem 3D}}}
            
            \small 
            
            \begin{itemize}
                \item[\faAngleRight] SolidWorks, FreeCAD, Blender.
            \end{itemize}
            
            \normalsize
            
            \vspace*{-1mm}
            
            
            \section{\textbf{\faCalculator\hspace{1mm}\en[\textit{Hardware}]{Hardware}}}
            
            \small 

            \begin{itemize}
                \item[\faAngleRight] \en{Controllers and computers}\fr{Microcontrôleur et ordinateurs}\pt{Controladores e computadores}
                \begin{itemize}
                    \item  Nvidia Jetson, Arduino, Intel MCS-51, Atmel AVR, ARM Cortex, dSPACE, NI DAQ, IoT, RTOS, Windows, Linux.
                    % \item  Raspberry Pi 3, Arduino, Intel MCS-51, Atmel AVR, ARM Cortex, dSPACE, National Instruments DAQ.
                \end{itemize}
                \item[\faAngleRight] \en{Sensors}\fr{Capteurs}\pt{Sensores}
                \begin{itemize}
                    \item  LIDARs,\fr{ Caméras}\pt{ Câmeras} RGBD\en{ Cameras}, IMU, GNSS.
                \end{itemize}
                \item[\faAngleRight] \en{Others}\fr{Autres}\pt{Outros}
                \begin{itemize}
                    \item \en{Mechatronic prototyping, PCB manufacturing, Instrumentation, Communication Protocols}\fr{Prototypage mécatronique, fabrication de circuits imprimés, Instrumentation, pProtocoles de Communication}\pt{Prototipagem Mecatrônica, Fabricação de PCBs, Instrumentação, Protocolos de Comunicação}.
                \end{itemize}
            \end{itemize}
            
            
            \normalsize
            
            \vspace*{-1mm}
            \section{\textbf{\faGlobeAmericas\hspace{1mm}\en{Languages}\fr{Langues}\pt{Línguas}}}
            \vspace*{-3mm}
            \small
            \begin{center}
                \begin{tabular}{p{0.4\textwidth} p{0.6\textwidth}} 
                    \en{Portuguese}\fr{Portugais}\pt{Português}&\pt[Native]{Nativa}\\
                    \en{English}\fr{Anglais}\pt{Inglês}&\en{Fluent}\fr{Courant}\pt{Fluente}\\
                    \en{French}\fr{Français}\pt{Francês}&\en{Fluent}\fr{Courant}\pt{Fluente}\\
                \end{tabular}
            \end{center}
            \vspace*{-8mm}
            % \section{\textbf{\faChartPie\hspace{1mm}\en[\textit{Softskills}]{Softskills}}}
            % \small \en[\textit{Problem-Solving}]{Problem-Solving}, \en{adaptability e flexibility, driven to learn, autonomy, pro-activity. Familiar with Agile methodologies}\fr{adaptation et flexibilité, volonté d’apprendre, autonomie, proactivité. Habituée avec des methodologies Agile}\pt{adaptação e flexibilidade, vontade de aprender, autonomia, proatividade. Familiarizada com metodologias ágeis}.\normalsize
            
            
            % \section{\textbf{\faCloudsmith\hspace{1mm}\textit{Personal Interests}}}
            % \small \trig Literature, Philosophy, Digital Art, Game development, Chess.  \textbf{Approaches:} Biomimetics, Open Source, IoT.
            
            % \vspace*{-5mm}
            
            % \normalsize
        
        \end{tcolorbox}
    \end{minipage}
    \begin{minipage}[t]{0.68\textwidth}
      \vspace*{-0.2cm}
            
            % \begin{tcolorbox}[grow to right by=0.1cm,grow to left by=0.2cm,left = 0mm,right = 0mm,colback=white,colframe=white,arc=2mm]

            %         \textbf{Control and Automation (Mechatronics) engineering student searching for remote job after finishing her double degree in \france.}
            %     \end{tcolorbox}

            
            % \vspace*{-2mm}
            % \normalsize
            % \begin{tcolorbox}[colframe=bluegrey,colback=white,grow to right by=0.1cm,grow to left by=0.2cm,boxrule=2mm,bottomrule=5mm,sharp corners=south,title=\textbf{\faBullseye\hspace{1.5mm}Motivation}]
            %     \centering

            %     \textbf{Control and Automation (Mechatronics) engineering student searching for remote job after finishing her double degree in \france.}
                
            %     %\textbf{Mobilité: \france}
                
                
            % \normalsize
            % \end{tcolorbox}
            % %sharp corners=south
        
            % \vspace*{-0.5cm}
        \begin{tcolorbox}[grow to right by=0.1cm,grow to left by=0.2cm,left = 0mm,right = 0mm,colback=white,colframe=white,arc=2mm]
            
            
            \setlength{\hintscolumnwidth}{2cm} 
            
            %\begin{tcolorbox}[blanker,colback=bluegrey,left=5mm,right=5mm, borderline vertical={4pt}{0pt}{darkdarkgrey}]
                
            %\end{tcolorbox}

            %Python, C,\CPP, instrumentation, traitement du signal, techniques d’estimation et fusion des mesures, robotique (sous-)marine, ROS 1/2, planification de trajectoire, planification automatique de tâches, optimisation numérique
            
            %\CPP moderne, modélisation des robots et concéption mécanique
            %\CPP, ROS, ROS 2, modélisation des robots
            %\CPP moderne, robotique mobile planification
            %ROS, \CPP moderne, modélisation et commande de robots, optimisation
            %ROS, \CPP, instrumentation, prototypage, électronique pour traitement du signal
            %Instrumentation, design et le prototypage de systèmes électroniques
            \vspace*{-5mm}
            \normalsize
            \section{\textbf{\faGraduationCap\hspace{1mm}\en{Academic}\fr{Formation}\pt{Formação Acadêmica}}}
            \vspace*{-3mm}
            
            ((% if use_double_degree_tcolorbox %))
            \begin{tcolorbox}[breakable,colback=white, left = -2mm,right = -2mm, boxrule=1mm, title=\textbf{\hspace{2mm} \faGraduationCap \hspace{1mm}
            \en{Double Degree}
            \fr{Double Diplôme}
            \pt{Duplo Diploma}}]
            \vspace{-1mm}
            ((% endif %))
            
                % \academicitem{\september\text{ }2020 \newline \hspace*{-1mm}\september\text{ }2022}
            % {ECN - École Centrale de Nantes}
            % {Nantes, \france}
            % {\en{\textbf{Engineering degree with specialization in Robotics and R\&D (Master's degree)}. \textit{Grand École d'Ingénieurs}, 4$^\text{th}$ place in France according to \textit{l'Étudiant} in 2021.}\fr{\textbf{Diplôme d'ingénieur.e avec spécialisation en Robotique et Recherche et Développement \textit{(Master's degree)}}. Grand École d'Ingénieurs, 4$^\text{ème}$ place en France selon le classement de l'Étudiant en 2021.}\pt{\textbf{Diploma de Engenharia com especialização em Robótica e em Pesquisa e Desenvolvimento \textit{(Master's degree)}}. \textit{Grand École d'Ingénieurs}, 4$^\circ$ lugar na França de acordo com o \textit{l'Étudiant} em 2021.}}
              % #1: id, #2: institution, #3: logo, #4: location, #5: start, #6: end, #7: desc, #8: skills
              \academicitemtimeline
                  {ecn}
                  {ECN - École Centrale de Nantes} % Institution
                  {((parent_prefix))pictures/ecn_logo.jpg} % Logo
                  {Nantes, France} % Location
                  {Sep 2020} % Start
                  {Sep 2022}  % End
                  {\textbf{Engineering degree with specialization in Robotics and R\&D (Master's degree)}. \textit{Grand École d'Ingénieurs}, 4$^\text{th}$ place in France according to \textit{l'Étudiant} in 2021.} % Description
                  {Robotics Modeling, Non-Linear Control, Electronics, C++, ROS 2, Computer Vision.} % Skills
              % \academicitem{\february\text{ }2017 \newline \december\text{ }2020 \newline\vfill \november\text{ }2022 \september\text{ }2023}
              % {UFRGS - Universidade Federal do Rio Grande do Sul\newline}{Porto Alegre, \brazil}
                  % {\en{\textbf{Control and Automation Engineering}. Among the best universities in \brazil\text{ }and twice the best control and automation/mechatronics engineering program according to ENADE.}\fr{\textbf{Ingénierie Mécatronique}. Parmi les meilleurs universités brésiliennes et le meilleur cours d'Ingénierie Mécatronique du Brésil deux fois consécutives selon ENADE.}\pt{\textbf{Engenharia de Controle e Automação}. Entre as melhores universidades brasileiras e melhor curso de Engenharia de Controle e Automação do Brasil duas vezes consecutivas segundo o ENADE.}}
              \academicitemtimeline
                  {ecaufrgs}
                  {UFRGS - Universidade Federal do Rio Grande do Sul} % Institution
                  {((parent_prefix))pictures/ufrgs.png} % Logo
                  {Porto Alegre, Brazil} % Location
                  {Feb 2017} % Start
                  {Nov 2023}  % End
                  {\en{\textbf{Control and Automation Engineering (Mechatronics)}. Among the best universities in \brazil\text{ }and twice the best control and automation/mechatronics engineering program according to ENADE.}\fr{\textbf{Ingénierie Mécatronique}. Parmi les meilleurs universités brésiliennes et le meilleur cours d'Ingénierie Mécatronique du Brésil deux fois consécutives selon ENADE.}\pt{\textbf{Engenharia de Controle e Automação}. Entre as melhores universidades brasileiras e melhor curso de Engenharia de Controle e Automação do Brasil duas vezes consecutivas segundo o ENADE.}} % Description
                  {Control Theory, Automation, Electronics, Electric Machines, Solid Mechanics, Mechanisms, Mathematical Modeling, C.} % Skills
                  % {Control, Electronics, Electric Machines, Mathematical Modeling, C.} % Skills
                
            \vspace*{-4mm}
            ((% if use_double_degree_tcolorbox %))
            \end{tcolorbox}
            \vspace*{-4mm}
            ((% else %))
            \vspace*{-1mm}
            ((% endif %))
            
%--------------Expérience------------------------            
          \section{\textbf{\faSuitcase\hspace{1mm}\en{Professional Experience}\fr{Expérience Professionnelle}\pt{Experiência Profissional}}}

            \vspace*{-6mm}
            \begin{center}
          ((% for job in experiences %))
              \jobitemtimeline
                  {((job.id))} 
                  {((job.company))} 
                  {((parent_prefix))((job.logo_path))} 
                  {((job.location))} 
                  {((job.role))} 
                  {((job.start_date))} 
                  {((job.end_date))}  
                  {((job.intro))((job.description))} 
                  {(( job.skills | join(', ') ))}
          ((% endfor %))
          \end{center}
          ((% set ns = namespace(prev_id=none) %))
          ((% for job in experiences %))
              ((% if ns.prev_id is not none %))
                  \connectjobs{((ns.prev_id))}{((job.id))}
              ((% endif %))
              ((% set ns.prev_id = job.id %))
          ((% endfor %))

          \vspace*{-8mm}
            
        \end{tcolorbox}
    \end{minipage}
    
    \vspace*{-2mm}
\end{tcolorbox}


\end{document}
